\section{Outlook: would a larger GPU change the answer?}
\label{sec:h100}

Partly, and the claim should be scoped accordingly. The dominant cause,
Sect.~\ref{sec:dilemma}, is a \emph{capacity} limit, and datacenter GPUs relax it
substantially (Table~\ref{tab:h100}). With 228\,KB of shared memory per SM one could hold
${\approx}390$ zones on-chip, so R2's failure mode largely disappears. Because the workload is
latency-bound with millions of \emph{independent} zones, absolute throughput should also
scale with GPU size even at unchanged utilisation; scaling crudely by lanes and bandwidth
suggests $50$--$70$\,M succ/s against our CPU's $17$\,M.

\begin{table}[t]
\centering
\caption{The dominant cause is substantially an Apple-Silicon capacity constraint.}
\label{tab:h100}
\begin{tabular}{lrr}
\toprule
& Apple M4 Pro & NVIDIA H100 \\
\midrule
Shared / threadgroup memory & 32\,KB per group & 228\,KB per SM \\
Memory bandwidth & 273\,GB/s (shared with CPU) & ${\approx}3.3$\,TB/s dedicated \\
Lanes & ${\approx}2\,000$ & ${\approx}16\,900$ \\
\bottomrule
\end{tabular}
\end{table}

What would \emph{not} improve: the serial $k$ loop, and the counting argument of
Sect.~\ref{sec:cooperation} (warps are also 32 lanes wide), so one-zone-per-thread remains the
right mapping --- merely a wider one. Coherency becomes \emph{worse} on a discrete GPU, since
PCIe transfers replace cache snooping unless the entire search is made GPU-resident, which in
turn requires a correct lock-free device-side hash table; our attempt at one produced
incorrect results because Metal Shading Language exposes only relaxed atomics, with no
acquire/release pairing for the publication step.

Two caveats for a fair follow-up: an H100 should be compared against a similarly priced CPU
(a 128-core server part, not 12 consumer cores), and GPU residency is a prerequisite rather
than an optimisation.
